
A strong correlation (r=0.80) exists between LLM truthfulness and adversarial robustness---two capabilities often studied independently---but the intuitive explanation involving calibration is not supported by the evidence.

This finding has implications for AI safety research: if models that avoid generating misinformation also resist adversarial attacks, understanding the underlying mechanism could inform the development of trustworthy systems.

\subsubsection{1.1 The Problem: Trust Benchmark Silos}

Large language models fail in multiple trust dimensions. A model may generate plausible but factually incorrect statements (truthfulness failures), or it may be misled by adversarially perturbed inputs that preserve semantic meaning (robustness failures). These failure modes are typically studied independently: TruthfulQA \citep{lin2022truthfulqa} evaluates whether models repeat common misconceptions, while AdvGLUE \citep{wang2022advglue} measures vulnerability to word-level adversarial perturbations.

This separation leaves open the question of whether truthfulness and robustness are fundamentally related capabilities or independent dimensions. If correlated, a common mechanism may underlie both; if independent, separate interventions would be required. No prior study has systematically correlated TruthfulQA and AdvGLUE scores across multiple model families.

\subsubsection{1.2 Approach and Key Finding}

This work conducts the first systematic correlation analysis between truthfulness (TruthfulQA MC1) and adversarial robustness (AdvGLUE average) across 14 decoder-only LLMs from four families (Pythia, Llama-2, Mistral, Falcon), ranging from 70M to 70B parameters. Using partial correlation with log-transformed model size as a covariate:

\textbf{Main Result:} TruthfulQA MC1 and AdvGLUE accuracy show a strong positive partial correlation (r=0.8028, p=0.000548) with a 95\% bootstrap confidence interval of [0.0816, 0.9685] that excludes zero.

The calibration hypothesis---that well-calibrated models accurately estimate uncertainty, enabling them to avoid false claims and detect anomalous inputs---is tested directly by measuring Expected Calibration Error (ECE) on MMLU.

\textbf{Mechanism Finding:} The calibration hypothesis is not supported. ECE shows near-zero correlation with TruthfulQA (r=-0.12, p=0.68) and AdvGLUE (r=-0.16, p=0.58). Furthermore, poorly-calibrated (high-ECE) models showed stronger truthfulness-robustness correlation (r=0.99) than well-calibrated models (r=0.65)---the opposite of what the calibration mechanism predicts (Fisher's z-test p=0.165).

\subsubsection{1.3 Contributions}

\begin{enumerate}
\item \textbf{First documented correlation:} TruthfulQA MC1 and AdvGLUE accuracy are strongly positively correlated (r=0.80) after controlling for model size.
\end{enumerate}

\begin{enumerate}
\item \textbf{Mechanism falsification:} Model calibration (ECE) does not explain this correlation, ruling out the most intuitive hypothesis.
\end{enumerate}

\begin{enumerate}
\item \textbf{Scope characterization:} The correlation is confirmed for base models (r=0.80, p=0.0005, N=8). Results for instruction-tuned models (N=6) are inconclusive due to limited sample size.
\end{enumerate}
